\section{关键组件详析}
\label{sec:components}

本节按照从核心到外设的顺序，依次分析六个关键组件：
基础模型（对应CPU）→ KV缓存（对应缓存层次）→ 上下文管理（对应虚拟内存）→ Agent运行时（对应操作系统）→ 工具总线（对应I/O）→ 多Agent协作（对应分布式系统）。
对于每个组件，我们都遵循四步结构：
(1) 经典计算机中的对应概念是什么；(2) 大模型系统中的对应物是什么；(3) 两者为什么相似（机制层面的解释）；(4) 本质差异是什么（不回避问题）。

%% ---------- 基础模型：CPU ----------
\subsection{基础模型：概率式模型计算核}
\label{sec:model}

\subsubsection{经典计算机中的CPU}

在经典计算机中，CPU（中央处理器）是整个系统的执行核心。
它的职责可以用一句话概括：\textbf{按指令序列执行操作，将输入状态映射为输出状态}。
CPU内部包含算术逻辑单元（ALU）用于计算、寄存器文件用于暂存操作数、控制单元用于取指译码。
关键特征是\textbf{确定性}：执行\texttt{ADD R1, R2, R3}一百万次，只要输入寄存器的值相同，结果就一定相同。
CPU的行为由指令集架构（ISA）精确定义，例如RISC-V规范对每条指令的行为给出了数学级的描述\cite{riscv_spec_2026}。

\subsubsection{大模型系统中的对应物}

在大模型系统中，基础模型（如GPT-4、Claude、Llama）承担了类似CPU的``通用执行核心''角色。
给定输入上下文（包含系统提示、用户消息、工具结果与控制约束），模型通过前向计算产生下一步动作或输出\cite{vaswani2017attention,brown2020gpt3,llama2023}。

提示词、工具描述、Skill文件、项目记忆文件（如\texttt{AGENTS.md}或\texttt{CLAUDE.md}）可以被理解为一种\textbf{软指令接口}：它们不改变模型的物理权重（就像软件不改变CPU的电路），却显著改变模型的执行路径\cite{openai_codex_agents_md_2026,openai_codex_skills_2026,anthropic_claudecode_memory_2026}。
例如，一个精心编写的\texttt{CLAUDE.md}文件可以让同一个模型在Python项目中表现出``Python专家''的行为模式，在Rust项目中表现出``Rust专家''的行为模式——就像同一个CPU加载不同的程序就表现出不同的行为。

SGLang已经朝这个方向迈进：它把复杂的语言模型程序重新表述为结构化前端语言与高性能运行时的协同，类似于从汇编语言向高级语言的演化\cite{sglang2024,sglang_docs_2026}。

\subsubsection{为什么相似}

CPU和基础模型之间的相似性体现在三个层面：

\textbf{第一，通用性}
CPU是通用计算核心——同一个CPU可以运行文字处理器，也可以运行科学模拟程序。
基础模型也是通用推理核心——同一个模型可以写代码、分析数据、翻译语言、规划任务。
这种通用性使得两者都可以作为``硬件基础''，上层通过不同的``软件''（程序 vs.\ 提示词）来定制行为。

\textbf{第二，接口抽象}
CPU通过ISA向上层暴露操作接口（算术、跳转、访存）；模型通过提示模板和工具schema向上层暴露操作接口（推理、生成、工具调用）。
两者都遵循``接口稳定、实现可演进''的原则——x86架构存在了40多年，而底层微架构从单标量演进到了超标量乱序执行；GPT-4的API接口保持稳定，而底层模型版本持续更新。

\textbf{第三，微架构优化空间}
CPU有流水线、分支预测、超标量等微架构优化；模型有FlashAttention（IO感知的注意力计算\cite{dao2022flashattention}）、推测解码（用小模型预测大模型输出\cite{leviathan2023speculative}）、连续批处理等微架构优化。
两者都面临``接口不变、底层性能持续提升''的工程挑战。

\subsubsection{本质差异}

差异的核心在于\textbf{确定性 vs.\ 概率性}。
CPU执行形式化指令集，每条指令的语义精确定义、长期稳定、完全可复现。
模型处理的是自然语言、代码和结构化schema的混合输入，输出服从概率分布。
正如Mi等人所指出的，从计算机体系结构的演进中汲取灵感来设计Agent系统，需要同时承认概率式执行与确定性执行的根本差异\cite{mi2025llmagentscomputersystems}。

因此，对未来``智能ISA''的理解不应是把提示词硬编码为固定指令，而应是发展更稳定的接口层：工具schema、受控输出格式、可复用技能包、任务DSL以及对上下文的显式声明。
从体系结构观点看，这意味着未来研究不应只比较模型参数量，还应比较``模型可编程性''——即上层能在多大程度上稳定、可靠地控制模型行为。

%% ---------- KV缓存：缓存层次 ----------
\subsection{KV缓存：智能系统的缓存层次}
\label{sec:kv}

\subsubsection{经典计算机中的缓存层次}

在经典计算机中，缓存是解决``处理器快、内存慢''这一根本矛盾的关键机制\cite{drepper2007memory}。
现代CPU通常有三层缓存：L1（约1ns、64KB）→ L2（约4ns、512KB）→ L3（约12ns、数MB到数十MB）。
缓存利用\textbf{时间局部性}（temporal locality）：最近被访问的数据很可能再次被访问。
当CPU需要数据时，先查L1，未命中则查L2，再未命中则查L3，最终从主存加载。
缓存的管理涉及三个核心问题：\textit{放置策略}（数据存在哪里）、\textit{替换策略}（满了淘汰谁）、\textit{写策略}（修改如何同步）。

\subsubsection{大模型系统中的对应物}

在自回归推理中，生成每个新token时都需要用该token的Query与之前\textit{所有}token的Key/Value计算注意力。
如果每生成一个token都重算整个前缀的KV，计算量随序列长度呈$O(n^2)$增长。
KV缓存通过保存已计算的历史token的Key和Value向量，将每步注意力计算简化为$O(1)$的查表操作，使整体推理从$O(n^2)$降为$O(n)$。
这与经典缓存的目的高度一致：\textbf{牺牲容量（显存）换取延迟与带宽优势}\cite{drepper2007memory,kwon2023pagedattention}。

\textbf{KV缓存的层次结构正在形成：}
\begin{itemize}[nosep]
\item \textbf{L1——会话级KV缓存}：单次会话中的KV向量，延迟最低，容量受GPU显存限制
\item \textbf{L2——前缀/共享缓存}：跨请求共享的系统提示或代码库索引的KV块，类似于CPU中多进程共享的只读代码页\cite{vllm_prefix_2026,sglang_docs_2026,gim2024promptcache}
\item \textbf{L3——语义缓存}：相同问题类下可复用的工具调用结果、相同仓库结构下可复用的索引，属于``语义级别''的缓存命中
\end{itemize}

当vLLM引入PagedAttention，把KV组织成固定大小block并通过block table映射到非连续物理空间时，这种类比已经不只是启发式，而是直接借用了虚拟内存的页式管理思想\cite{kwon2023pagedattention}。

\subsubsection{为什么相似}

KV缓存与CPU缓存的相似性体现在工程机制的层面：

\textbf{第一，都利用时间局部性}
CPU缓存保存最近访问的缓存行，因为``最近用过的数据很快会再用''；KV缓存保存最近计算的KV向量，因为``生成下一个token时必须 attends to 所有之前的token''。
两者都把``需要反复使用的数据''保存在高速存储中，避免昂贵的重复计算或访存。

\textbf{第二，都面临容量压力与替换决策}
L1缓存容量有限，需要LRU等替换策略决定淘汰哪些缓存行；KV缓存同样面临显存压力，H2O通过识别``重击者''token（attention权重最大的token）实现了智能驱逐策略\cite{h2o2023}，类似于CPU缓存中识别``热数据''的优化。
StreamingLLM发现初始token获得的注意力不成比例地高（``attention sinks''现象），通过保留这些sink token实现了任意长度的流式推理\cite{xiao2024streamingllm}——这类似于CPU缓存中的``钉住''（pinning）策略。

\textbf{第三，都存在带宽瓶颈}
Gholami等人指出，LLM推理在解码阶段是典型的内存带宽受限问题，而非计算受限\cite{gholami2024memorywall}。
Yuan等人用Roofline模型系统分析了这一现象\cite{yuan2024llminference}。
Li等人2024年的综述将KV缓存管理策略系统分为token级（决定哪些KV条目参与注意力）、模型级（架构优化如GQA/MQA）和系统级（内存管理与调度）三个层次\cite{likvcache2024survey}——这恰好对应了CPU缓存优化的多层次方法。

\textbf{第四，压缩是共同的优化方向}
KVQuant实现了8倍压缩\cite{hooper2024kvquant}；KIVI通过非对称2-bit量化在几乎无损的条件下大幅减少显存占用\cite{liu2024kivi}。
这类似于CPU缓存中通过数据压缩（如IBM MXT）增加有效容量的思路。

\subsubsection{本质差异}

\textbf{第一，命中判定标准不同}
CPU缓存命中是\textbf{布尔判断}——地址相等则命中，否则不命中。
KV缓存中，前缀缓存的确是精确匹配（前缀token序列相同则复用），但上层的语义缓存命中是\textbf{连续值比较}——语义相似度超过阈值即视为命中，不同阈值导致不同的精确率/召回率权衡。
这意味着传统缓存理论中的命中率模型不能直接适用于语义缓存。

\textbf{第二，失效机制不同}
CPU缓存失效由写入一致性协议（如MESI）保证——数据被修改时，相关缓存行被标记为失效。
KV缓存和语义缓存的``失效''远为复杂：工具版本变化、代码仓库提交变化、索引过期都会导致缓存命中但答案失真——即``语义失效''。
近期研究已开始探索学习式的驱逐策略\cite{lpc2025}和依赖感知的失效机制\cite{multilevelcache2025}来应对这一挑战。

\textbf{第三，容量-延迟权衡更为严峻}
CPU缓存层次中，L1→L2→L3→主存的延迟以纳秒级递增；在KV缓存中，从GPU显存→CPU内存→磁盘/网络的延迟跳变达到毫秒级，中间没有平滑过渡。
未来缓存管理需要把传统cache policy与provenance、TTL、哈希版本一起设计。

%% ---------- 上下文管理：虚拟内存 ----------
\subsection{上下文管理：语义版虚拟内存}
\label{sec:memory}

\subsubsection{经典计算机中的虚拟内存}

在经典计算机中，虚拟内存解决了一个核心矛盾：``程序需要的地址空间远大于物理内存容量''。
操作系统通过\textbf{分页机制}制造了``大空间''的幻觉：每个进程拥有独立的虚拟地址空间（如48位地址，256TB），但物理内存可能只有16GB。
当进程访问一个尚未映射到物理内存的虚拟页时，触发``页缺失''（page fault），操作系统从磁盘加载所需页面。
分页的核心优势是\textbf{按需加载}——只有实际被访问的数据才占用物理内存，且这个过程对应用程序完全透明。

虚拟内存管理的关键组件包括：页表（虚拟地址到物理地址的映射）、TLB（加速地址转换的缓存）、
换页策略（决定哪些页被换出到磁盘）和写回策略（脏页何时写回磁盘）\cite{ostep_2023}。

\subsubsection{大模型系统中的对应物}

上下文窗口之于大模型，正如主存之于进程：它是高带宽、低延迟、但容量有限的\textbf{工作区}。
Anthropic将context window直接定义为模型生成时可回看的working memory。

当前，上下文窗口容量正在快速扩展：Gemini 1.5 Pro支持高达1000万token的多模态上下文\cite{gemini2024}；
Ring Attention通过在多设备间环形分布序列实现了近乎无限的上下文长度\cite{liu2024ringattention}；
RoPE/YaRN/LongRoPE等位置编码方法旨在增加可寻址历史范围\cite{su2021roformer,yarn2024,longrope2024}。

但大量研究表明，``声明可支持的窗口长度''与``真正有效的工作集容量''不是一回事。
RULER、LongBench v2和LOFT对长上下文能力给出了更现实的压力测试\cite{ruler2024,longbenchv2_2024,loft2024}。
就像一台宣称有256TB虚拟地址空间的计算机，实际可用内存可能只有16GB，真正的性能取决于工作集（working set）是否常驻物理内存。

因此，未来更合理的设计不应是``把所有东西一次性塞进上下文''，而应像虚拟内存那样做\textbf{分层管理}。

\textbf{正在形成的分层方案：}
\begin{itemize}[nosep]
\item \textbf{热记忆（上下文窗口内）}：当前任务的完整信息，类似物理内存中的活跃页
\item \textbf{温记忆（摘要与索引）}：经压缩的历史上下文，类似换出但保留摘要信息的页
\item \textbf{冷记忆（外部知识库）}：持久化的向量库、文档库，类似磁盘上的数据
\end{itemize}

MemGPT把这种思想称为\textbf{虚拟上下文管理}，明确提出LLM作为操作系统的类比，将信息在上下文窗口（``主存''）和外部存储（``磁盘''）之间分页调度\cite{memgpt2023}。
LongMem则把长期记忆显式地从主干模型中解耦出来\cite{longmem2023}。
RAG在这里扮演的更像\textbf{外存页入机制}——当Agent需要``页缺失''的信息时，通过语义检索从知识库中加载相关内容到上下文窗口\cite{lewis2020rag}。
MemoryOS更进一步，直接将OS内存管理原则应用于Agent记忆系统\cite{memoryos2025}。
MemOS提出了完整的记忆操作系统架构\cite{memos2025}。
HiAgent受人类问题解决过程启发，用子目标作为记忆块来层次化管理Agent工作记忆\cite{hiagent2024}。
A-MEM提出让Agent自主组织记忆结构，无需预定义schema\cite{amem2025}。

\subsubsection{为什么相似}

上下文管理与虚拟内存的相似性体现在\textbf{机制层面}：

\textbf{第一，都制造``大空间''幻觉}
虚拟内存让每个进程以为自己拥有连续的大地址空间，实际物理内存远小于此。
虚拟上下文管理让Agent以为自己拥有无限记忆，实际上下文窗口只有128K-1M token。
两者都通过\textbf{按需加载}实现幻觉：只有当前需要的信息才占据``物理内存''（上下文窗口），其余信息存储在``磁盘''（外部知识库）中。

\textbf{第二，都面临``地址转换''问题}
虚拟内存通过页表将虚拟地址映射到物理地址；虚拟上下文通过语义检索将``信息需求''映射到具体的知识条目。
只不过前者是精确匹配（页号→物理页号），后者是近似匹配（查询意图→Top-K相关文档）。

\textbf{第三，都需要``换页策略''}
当物理内存不足时，OS必须决定哪些页被换出（LRU、Clock等策略）；当上下文窗口快满时，Agent必须决定哪些信息被移出（摘要压缩、选择性遗忘）。
这正是``上下文编译器''的核心功能——系统必须决定哪些状态以原文加载，哪些以摘要加载，哪些只保留索引，哪些应被丢弃\cite{contextcompiler2025}。

\subsubsection{本质差异}

\textbf{关键差异：虚拟内存是\textbf{无损}的，上下文管理是\textbf{有损}的}
OS的分页机制对应用程序完全透明——页换入换出不会改变任何一个字节。
而Agent的上下文管理涉及\textbf{语义压缩}——摘要必然丢失细节，选择必然遗漏信息。
``地址转换''也不是精确的：向量相似度搜索返回的是``最相关的''结果，而非``完全匹配的''结果。

这意味着语义虚拟内存永远不是无损页式系统。
它更接近``带语义摘要和证据回填的分页内存''——信息被压缩后可能失真，但可以通过溯源（provenance）链回到原始数据验证。
JetBrains的研究指出，高效上下文管理的关键不在于``塞入更多token''，而在于``减少噪声、保留信号''\cite{jetbrains2025context}——这进一步验证了有损压缩的必要性。

%% ---------- Agent运行时：OS ----------
\subsection{Agent运行时：智能操作系统}
\label{sec:agent}

\subsubsection{经典计算机中的操作系统}

在经典计算机中，操作系统是``管理硬件资源、隔离应用程序、提供统一服务接口''的系统软件\cite{ostep_2023}。
它的核心职责可以用三个抽象概括：
\textbf{虚拟化}（让每个进程以为独占CPU和内存）、\textbf{并发}（让多个进程安全地同时运行）、\textbf{持久化}（可靠地存储和管理文件）。

具体来说，OS要做的事情包括：响应系统调用（如文件读写）、调度进程（哪个先运行、运行多久）、隔离进程（一个崩溃不影响另一个）、管理内存（分配、回收、换页）、控制权限（谁能访问什么资源）、处理I/O（与外部设备交互）、记录日志（审计与故障恢复）。

\subsubsection{大模型系统中的对应物}

如果模型是CPU，那么Agent框架越来越像OS。
Agent运行时要做的事情包括：
\begin{itemize}[nosep]
\item \textbf{解析目标}：将用户请求分解为可执行的子任务（类似程序加载与解析）
\item \textbf{决定是否规划}：简单请求直接执行，复杂请求需要多步规划（类似短进程快速路径 vs.\ 长进程调度）
\item \textbf{装入上下文}：加载系统提示、项目记忆、历史对话到上下文窗口（类似进程的内存映像加载）
\item \textbf{执行工具调用}：通过MCP等协议调用外部工具（类似系统调用）
\item \textbf{审批危险动作}：对有副作用的操作（文件写入、代码执行）进行权限检查（类似capability-based访问控制）
\item \textbf{并发启动子Agent}：将子任务分派给独立的子Agent执行（类似fork新进程）
\item \textbf{整合中间结果}：汇总子Agent的输出，形成最终答案（类似进程间通信与结果合并）
\item \textbf{状态持久化}：将重要状态写入记忆文件或知识库（类似文件系统写入）
\item \textbf{失败恢复}：出错时重试或回滚，保证任务一致性（类似事务恢复）
\end{itemize}

这一类比的学术化表达已经出现。
Ge等人提出``LLM as OS, Agents as Apps''的愿景\cite{ge2023llmos}。
Mei等人提出AIOS，构建了完整的LLM Agent操作系统内核架构\cite{mei2024aios}。
Mi等人系统地将计算机体系结构的演进经验应用于Agent系统设计\cite{mi2025llmagentscomputersystems}。
Karpathy在2023年指出``LLM不是聊天机器人，而是新操作系统的内核进程''%
\cite{karpathy2023llmos}——文件系统变成向量数据库，浏览器变成互联网搜索。
这一愿景正在成为经济现实：一家Augment Code的企业客户将其CTO预估需要4--8个月的项目在2周内完成\cite{anthropic2026agenticreport}，证明了有效的L3上下文管理与L5编排相结合可以将开发者入职和项目交付周期从周级别压缩到天级别。

工业实践也在朝这个方向演化：
Codex提供subagents、skills、sandbox等控制面\cite{openai_codex_overview_2026}；
Claude Code提供hooks、MCP、memory文件与只读审批\cite{anthropic_claudecode_overview_2026}。

\subsubsection{为什么相似}

Agent运行时与操作系统的相似性是所有类比中最为深刻的，因为它体现在\textbf{架构职责}的层面：

\textbf{第一，都是资源管理者}
OS管理CPU时间片、内存页、磁盘空间、网络带宽；Agent运行时管理模型推理配额、上下文窗口空间、工具调用频率、子Agent数量。
两者都需要回答：``有限的资源如何分配给竞争的需求？''

\textbf{第二，都提供隔离与安全}
OS通过地址空间隔离进程，一个进程崩溃不影响其他进程\cite{ostep_2023}；Agent运行时通过沙箱隔离子Agent，一个子Agent的错误操作不影响主任务\cite{openai_codex_sandbox_2026,anthropic_claudecode_sandbox_2026}。
两者都面临权限控制问题——``谁能做什么''。

\textbf{第三，都定义标准接口}
OS通过系统调用（POSIX接口）标准化应用程序与硬件的交互；Agent运行时通过MCP等协议标准化模型与外部工具的交互\cite{mcp_intro_2026}。

\subsubsection{本质差异}

\textbf{第一，Agent``进程''的行为不可预测}
传统OS调度的进程，其行为由确定性代码定义，操作系统可以精确预测执行时间和资源需求。
Agent的行为由概率式推理决定，同样的任务可能走上完全不同的执行路径——一次成功，另一次可能陷入死循环。
这给调度和资源管理带来了根本性挑战。

\textbf{第二，共享状态更脆弱}
多进程共享内存时，一致性由硬件缓存协议（如MESI）和软件锁机制保证——精确、可验证。
多Agent共享语义状态（如``对同一份文档的理解''）时，一致性只能通过自然语言沟通或外部同步机制维护——模糊、易出错。

\textbf{第三，接口定义不严格}
POSIX接口的每个系统调用都有精确定义的语义、参数类型、错误码和边界条件。
Agent的工具调用接口虽然可以用JSON Schema描述参数，但工具的``语义行为''（什么情况下返回什么结果）远未达到系统调用级别的精确度。

\textbf{这类比最有价值之处}，是让许多Agent问题可以被重新表达为OS问题：
\begin{itemize}[nosep]
\item 多Agent协作像多线程还是多进程？\cite{wu2023autogen}
\item 工具调用算不算系统调用？
\item MCP server是否相当于设备驱动或总线协议？\cite{mcp_intro_2026}
\item 危险命令为何需要capability-based权限？\cite{debenedetti2025camel}
\item 长任务如何通过日志和中断恢复？\cite{openai_codex_approvals_2026}
\end{itemize}

\textbf{安全层面的深刻类比}
在安全层面，这一类比尤为深刻。
CaMeL将操作系统能力安全原则应用于LLM Agent，创建了独立于LLM的保护层\cite{debenedetti2025camel}。
IronClaw明确提出要像操作系统一样保护AI Agent\cite{ironclaw2025}。
系统安全研究者开始用``agentic computing''这一新术语来定义Agent安全的基础问题\cite{agenticsecurity2025}。
这些工作说明，\textbf{安全边界首先是``运行时结构问题''，其次才是提示词问题}——正如OS安全首先取决于权限模型和隔离机制，其次才取决于应用程序的编码质量。

%% ---------- 工具总线与I/O ----------
\subsection{工具总线与Agent互联}
\label{sec:io}

\subsubsection{经典计算机中的I/O系统}

在经典计算机中，I/O系统负责连接CPU/内存与外部世界。
它的演化历史极具启发性：
\textbf{早期阶段}（1960s-1970s），每个外设需要专用接口和驱动——显卡有自己的总线，硬盘有自己的控制器，打印机有自己的并行口。
\textbf{标准化阶段}（1990s-2000s），PCI/PCIe统一了主板级总线，USB统一了外设连接——硬件厂商只需遵循统一协议，设备即可即插即用。
\textbf{网络化阶段}（2000s至今），TCP/IP协议栈统一了网络通信——任何设备都可以通过标准协议互联。

这一历史的核心教训是：\textbf{标准化总线协议是生态系统繁荣的前提}。

\subsubsection{大模型系统中的对应物}

在模型原生计算系统中，I/O层面正在经历类似从``点对点''到``标准化总线''的转变。

\textbf{早期阶段（方法层面）：}
ReAct将reasoning与acting交织起来\cite{yao2023react}；Toolformer让模型学会何时调用API\cite{schick2023toolformer}；HuggingGPT把LLM放到``controller''位置来调度异构模型\cite{hugginggpt2023}；Gorilla聚焦API调用的准确性与文档更新适应性\cite{patil2024gorilla}。
这些工作相当于早期计算机中``直接操作外设''的阶段：每个工具都需要单独的接口和调用约定，缺乏统一的协议标准。

\textbf{标准化阶段（协议层面）：}
MCP（Model Context Protocol）将连接AI应用与外部数据源、工具、工作流标准化\cite{mcp_intro_2026}，扮演类似PCIe总线或USB在传统计算机中的角色——它定义了统一的发现、连接、调用和错误处理机制，使新工具可以``即插即用''。
Google的A2A（Agent-to-Agent）协议则定位于agent-to-agent的互联标准\cite{google_a2a_2025}，支持capability discovery、task lifecycle和多模态协同，扮演类似网络协议栈中TCP/IP的角色——它使不同供应商开发的Agent可以互相发现和协作。

把两者合起来看，MCP是model到系统的\textbf{纵向总线}（Agent→工具），A2A是agent到agent的\textbf{横向互联}（Agent→Agent）。
正如一篇综述所指出的，MCP、A2A与ACP等协议正在形成Agent互操作的基础设施层\cite{agentprotocols2025}。

\subsubsection{为什么相似}

\textbf{第一，都解决``接口多样性''问题}
PCIe统一了主板级设备接口之前，每个设备需要专用驱动；MCP统一了工具接口之前，每个工具需要专门的API适配代码。
两者都通过``定义标准协议、隐藏实现差异''来降低系统集成成本。

\textbf{第二，都有分层协议栈}
传统I/O有物理层（电气信号）→ 链路层（PCIe事务层）→ 协议层（NVMe/SCSI）→ 应用层（文件系统）。
Agent I/O也有类似分层：传输层（HTTP/SSE）→ 协议层（MCP/A2A消息格式）→ 语义层（工具schema/Agent capability描述）→ 应用层（具体工具调用或Agent协作）。

\textbf{第三，都面临发现与协商问题}
USB设备插入后需要``枚举''——主机询问设备类型、能力、驱动需求；MCP服务器连接后也需要``capability discovery''——Agent查询服务器支持哪些工具、每个工具的参数schema和副作用描述。

\subsubsection{本质差异}

\textbf{第一，调用的确定性不同}
传统系统调用（如\texttt{read(fd, buffer, size)}）的行为完全可预测——参数相同、结果相同。
工具调用的结果可能是非确定的（搜索引擎返回的结果随时间变化）、异步的（长时间运行的任务需要回调）且带副作用的（发送邮件、执行代码）。

\textbf{第二，错误处理的复杂度不同}
PCIe设备要么正常响应，要么返回明确的错误码；MCP工具可能返回语义模糊的错误（``结果不太确定''），甚至返回看似正确但实际有误的结果。
这要求Agent I/O层引入\textbf{验证和冗余}机制——对关键操作进行结果校验，类似于分布式系统中的quorum读取。

%% ---------- 多Agent与分布式 ----------
\subsection{多Agent协作：语义版分布式系统}
\label{sec:distributed}

\subsubsection{经典计算机中的分布式系统}

当单台计算机的计算能力不足以完成任务时，需要多台计算机通过网络协作，这就是分布式系统。
分布式系统的核心挑战包括：
\textbf{任务分解与调度}（把大任务拆成小任务，分配到不同节点）、
\textbf{状态同步与一致性}（多个节点如何就共享状态达成一致）、
\textbf{故障检测与恢复}（某个节点崩溃后系统如何继续运行）、
\textbf{死锁避免}（多个节点互相等待如何打破僵局）、
\textbf{负载均衡}（如何让所有节点都保持忙碌而非某些节点过载而其他空闲）\cite{hennessy2017quantitative}。

\subsubsection{大模型系统中的对应物}

当多个Agent协作完成复杂任务时，系统呈现出分布式计算的特征。

AutoGen把多agent会话作为通用基础设施\cite{wu2023autogen}——多个Agent像分布式节点一样通过消息传递协作。
MetaGPT用assembly line和SOP将多agent协作显式流程化\cite{metagpt2023}——类似于工业流水线式的任务分解，每个Agent负责一个环节（需求分析→设计→编码→测试）。
Internet of Agents从``互联网''而非单机系统来想象agent生态\cite{ioa2024}——强调agent integration protocol、动态组队和分布式环境。

从体系结构视角看，存在清晰的对应关系：
单Agent对应单进程，多Agent协作对应多线程或多进程，跨节点的Agent集群对应分布式系统。
由此引入经典的分布式系统问题：任务分解与调度、状态同步与一致性、故障检测与恢复、死锁避免、负载均衡。

\subsubsection{为什么相似}

\textbf{第一，都面临``通信开销 vs.\ 并行收益''的权衡}
分布式系统中，Amdahl定律告诉我们并行加速比受限于串行部分；多Agent系统中，定律III（Agent加速比定律）预测类似的行为——当编排效率$E$不够高时，单纯增加Agent数量收益递减。
AutoGen的多Agent实验验证了这一预测：Agent数量从2增加到8时，任务完成时间的减少幅度逐渐收窄\cite{wu2023autogen}。
工业案例进一步表明，协调拓扑对$E$有显著影响：Fountain采用``中心辐射''拓扑（一个中心Copilot协调筛选、文档生成、情感分析三个子Agent）实现了50\%更快的候选人筛选速度和72小时内的全面配置\cite{anthropic2026agenticreport}；MetaGPT采用流水线拓扑（SOP驱动的顺序协作）；AutoGen采用对等拓扑（对话式协调）。
不同的拓扑在不同的任务场景下产生不同的$E$值，这暗示定律III中的$E$不是固定的标量，而是与协调结构相关的设计变量。

\textbf{第二，都需要协调机制}
分布式系统用共识协议（如Raft\cite{raft2014}）让多个节点就决策达成一致；多Agent系统需要``辩论''``投票''或``层级裁决''等机制让多个Agent就结论达成一致。

\textbf{第三，都面临一致性与可用性的权衡}
CAP定理指出分布式系统不能同时满足一致性、可用性和分区容错性\cite{cap2002}；多Agent系统中，Agent之间的``语义一致性''（对同一事实的理解是否相同）与``响应速度''（是否等所有Agent都完成推理再给出结论）之间存在类似的张力。

\subsubsection{本质差异}

\textbf{第一，通信带宽与精度的根本差异}
分布式节点之间通过结构化协议通信（如gRPC的protobuf），消息语义精确无损；Agent之间通过自然语言或半结构化消息通信，信息在传递过程中可能被误解、遗漏或曲解——``语义噪声''是Agent分布式系统的独特挑战。

\textbf{第二，故障模式的差异}
分布式节点的故障通常是二元的（在线/离线）；Agent的``故障''可能是渐变的——它没有崩溃，但推理质量下降、偏离任务目标、产生幻觉。
检测``Agent是否在正确执行任务''远比检测``节点是否存活''困难。

\textbf{第三，组合涌现的不确定性}
分布式系统的行为是各节点行为的确定组合；多Agent系统的行为可能产生\textbf{涌现效应}——多个Agent各自正确执行，但整体产出不可预测的结果。
这对系统的可调试性和可问责性提出了全新要求。

%% ---------- 架构图与流程图 ----------

图~\ref{fig:arch}给出本文建议的模型原生计算分层架构示意。
图的每一层对应本文分析的一个组件层次：模型核心（CPU）→ KV缓存（缓存）→ 内存层（虚拟内存）→ Agent/OS层（操作系统）→ I/O与持久层（外部设备）→ 分布式底座（集群）。
注意Agent/OS层通过旁路直接访问内存层和I/O层，体现了双平面架构中确定性控制平面管理资源的职责。

\begin{figure}[htbp]
\centering
\resizebox{0.97\textwidth}{!}{%
\begin{tikzpicture}[
    node distance=4mm and 5mm,
    layer/.style={
        draw, rounded corners=4pt, minimum width=9.6cm,
        minimum height=1cm, font=\small, align=center,
        line width=0.7pt
    },
    subbox/.style={
        draw, rounded corners=3pt, minimum height=0.85cm,
        font=\footnotesize, align=center, line width=0.6pt
    },
    arrow/.style={-{Latex[length=2mm]}, thick, black!40},
    bypass/.style={
        -{Latex[length=2.8mm]}, line width=2pt,
        modelnative!85!black
    },
]

% ============================================================
% Background: Probabilistic Execution Plane  (L1--L3)
% ============================================================
\fill[classical!10, rounded corners=8pt]
    (-6.4,-6.5) rectangle (6.4,0.0);

% Background: Deterministic Control Plane  (L4--L5)
\fill[modelnative!10, rounded corners=8pt]
    (-6.4,0.6) rectangle (6.4,5.3);

% ============================================================
% L6 -- Application Layer (above both planes)
% ============================================================
\node[layer, fill=gray!18, draw=gray!55]
    (L6) at (0,6.2)
    {\textbf{L6 \, Application Layer}\\[-1pt]
     {\scriptsize User tasks \, agent apps \, workflow automation \, governance}};

% ============================================================
% L5 -- Orchestration Layer
% ============================================================
\node[layer, fill=modelnative!20, draw=modelnative!65]
    (L5) at (0,4.4)
    {\textbf{L5 \, Orchestration Layer}\\[-1pt]
     {\scriptsize Planning $\cdot$ Agent scheduling $\cdot$ Permission / Approval $\cdot$
      Failure recovery $\cdot$ State commitment}};

% ============================================================
% L4 -- Semantic Interface Layer
% ============================================================
\node[layer, fill=modelnative!20, draw=modelnative!65]
    (L4) at (0,2.8)
    {\textbf{L4 \, Semantic Interface Layer}\\[-1pt]
     {\scriptsize Tool schemas (Tool ABI) $\cdot$ Memory APIs $\cdot$ Trace APIs $\cdot$
      MCP / A2A protocols}};

% ---- Dashed interface line at L3/L4 boundary ----
\draw[dashed, thick, black!35, line width=0.9pt]
    (-6.4,0.3) -- (6.4,0.3);
\node[font=\tiny\itshape, text=black!45, anchor=west] at (-6.3,0.3)
    {L3--L4 interface contract};

% ============================================================
% L3 -- Context Management Layer
% ============================================================
\node[layer, fill=classical!20, draw=classical!60]
    (L3) at (0,-0.9)
    {\textbf{L3 \, Context Management Layer}\\[-1pt]
     {\scriptsize Hot / warm / cold tiering $\cdot$ Context compilation $\cdot$
      Summarization \& retrieval $\cdot$ Virtual context}};

% ============================================================
% L2 -- Inference Serving Layer (split into sub-components)
% ============================================================
\node[subbox, fill=classical!20, draw=classical!60,
      minimum width=4.4cm]
    (L2core) at (-2.2,-2.7)
    {\textbf{Model Core}\\[-1pt]
     {\scriptsize Foundation model / Inference strategy}};

\node[subbox, fill=classical!20, draw=classical!60,
      minimum width=2.3cm]
    (L2sess) at (1.5,-2.35)
    {\textbf{Session KV}\\[-1pt]
     {\scriptsize Per-request cache}};

\node[subbox, fill=classical!20, draw=classical!60,
      minimum width=2.3cm]
    (L2pref) at (1.5,-3.05)
    {\textbf{Prefix / Shared KV}\\[-1pt]
     {\scriptsize Cross-request reuse}};

% thin bracket grouping L2 sub-boxes
\draw[classical!50, line width=0.5pt]
    (3.0,-2.1) -- (3.25,-2.1) -- (3.25,-3.3) -- (3.0,-3.3);
\node[font=\tiny\bfseries, text=classical!70!black, anchor=west]
    at (3.35,-2.7) {L2};

% ============================================================
% L1 -- Physical Execution Layer
% ============================================================
\node[layer, fill=classical!20, draw=classical!60]
    (L1) at (0,-4.7)
    {\textbf{L1 \, Physical Execution Layer}\\[-1pt]
     {\scriptsize GPU / TPU / ASIC \, matrix ops $\cdot$ Attention kernels $\cdot$
      Sampling \& decoding}};

% ============================================================
% L0 -- Distributed Substrate (below both planes)
% ============================================================
\node[layer, fill=gray!18, draw=gray!55]
    (L0) at (0,-6.0)
    {\textbf{Distributed Substrate}\\[-1pt]
     {\scriptsize Continuous batching $\cdot$ Prefill / decode decoupling $\cdot$
      Parallel routing $\cdot$ Replication $\cdot$ Consistency}};

% ============================================================
% Plane labels (inside backgrounds, top-left)
% ============================================================
\node[font=\small\bfseries, text=classical!80!black, anchor=north west,
      fill=classical!6, fill opacity=0.7, text opacity=1,
      rounded corners=2pt, inner sep=2.5pt]
    at (-6.2,-0.1)
    {Probabilistic Execution Plane};

\node[font=\small\bfseries, text=modelnative!80!black, anchor=north west,
      fill=modelnative!6, fill opacity=0.7, text opacity=1,
      rounded corners=2pt, inner sep=2.5pt]
    at (-6.2,5.1)
    {Deterministic Control Plane};

% ============================================================
% Vertical arrows (inter-layer data / service flow)
% ============================================================
\draw[arrow] (L6)   -- (L5);
\draw[arrow] (L5)   -- (L4);
\draw[arrow] (L4)   -- (L3);
\draw[arrow] (L3)   -- (L2core);
\draw[arrow] (L2core) -- (L1);
\draw[arrow] (L1)   -- (L0);
\draw[arrow] (L2sess)  -- (L2core);
\draw[arrow] (L2pref)  -- (L2core);

% ============================================================
% BYPASS ARROWS -- the central visual feature
% ============================================================
% LEFT bypass: L5 down to L2 (KV cache management)
\draw[bypass]
    ([xshift=-1mm]L5.west) -- ++(-1.4,0)
    coordinate (bl1) -- ++(0,-5.9)
    coordinate (bl2) -- ++(1.4,0)
    ([xshift=-1mm]L2core.west);
\node[font=\scriptsize\bfseries, text=modelnative!85!black,
      rotate=90, anchor=south]
    at ([xshift=-8mm]bl2) {Direct KV cache control};

% RIGHT bypass: L4 down to L3 (context management) and L2
\draw[bypass]
    ([xshift=1mm]L4.east) -- ++(1.4,0)
    coordinate (br1) -- ++(0,-4.2)
    coordinate (br2) -- ++(-1.4,0)
    ([xshift=1mm]L3.east);

\draw[bypass]
    ([xshift=1mm]L5.east) -- ++(2.2,0)
    coordinate (br3) -- ++(0,-7.7)
    coordinate (br4) -- ++(-1.8,0)
    ([xshift=1mm]L1.east);
\node[font=\scriptsize\bfseries, text=modelnative!85!black,
      rotate=-90, anchor=south]
    at ([xshift=8mm]br4) {Direct inference control};

\end{tikzpicture}}
\caption{ICA layered architecture for model-native computing.
  \textbf{Top to bottom:} L6 Application, L5 Orchestration, L4 Semantic Interface
  (Deterministic Control Plane, orange), dashed L3--L4 interface boundary,
  L3 Context Management, L2 Inference Serving (Model Core + Session KV +
  Prefix/Shared KV), L1 Physical Execution (Probabilistic Execution Plane, blue),
  and Distributed Substrate.
  Thick orange bypass arrows show the Deterministic Control Plane reaching
  through the interface boundary to directly manage KV caches (L2) and
  hardware execution (L1).}
\label{fig:arch}
\end{figure}


图~\ref{fig:flow}给出一个典型请求的控制流。
该流程体现了双平面架构的交互方式：概率平面负责生成和理解，确定性平面负责调度、审批和状态管理。

\begin{figure}[htbp]
\centering
\resizebox{0.95\textwidth}{!}{%
\begin{tikzpicture}[
    % ---- Styles ----
    procB/.style={draw=classical!60, rounded corners=4pt,
        minimum width=6.2cm, minimum height=9mm, align=center,
        fill=classical!15, font=\small, line width=0.65pt},
    procO/.style={draw=modelnative!60, rounded corners=4pt,
        minimum width=6.2cm, minimum height=9mm, align=center,
        fill=modelnative!15, font=\small, line width=0.65pt},
    procN/.style={draw=gray!50, rounded corners=4pt,
        minimum height=9mm, align=center,
        fill=gray!10, font=\small, line width=0.65pt},
    dec/.style={draw, diamond, aspect=2.6, align=center,
        inner sep=2pt, font=\small, line width=0.65pt},
    arrB/.style={-{Latex[length=2.2mm]}, thick, classical!75},
    arrO/.style={-{Latex[length=2.2mm]}, thick, modelnative!75},
    arrG/.style={-{Latex[length=2.2mm]}, thick, gray!55},
    hdr/.style={font=\footnotesize\bfseries, inner sep=3pt,
        rounded corners=2pt, fill opacity=0.55, text opacity=1},
    node distance=5mm and 5mm,
]

% ============================================================
%  GEOMETRY
% ============================================================
\pgfmathsetmacro{\cxL}{-3.8}   % left-lane centre x
\pgfmathsetmacro{\cxR}{+3.8}   % right-lane centre x
\pgfmathsetmacro{\halfW}{3.5}  % half-width of each lane band

\def\yBot{-10.2}
\def\yTop{1.35}

% ============================================================
%  SWIM-LANE BACKGROUNDS  (blue left, orange right)
% ============================================================
\fill[classical!6, rounded corners=8pt]
    (\cxL-\halfW, \yBot-0.35) rectangle
    (\cxL+\halfW, \yTop+0.75);

\fill[modelnative!6, rounded corners=8pt]
    (\cxR-\halfW, \yBot-0.35) rectangle
    (\cxR+\halfW, \yTop+0.75);

% ---- Lane header pills ----
\node[hdr, text=classical!80!black, fill=classical!10]
    at (\cxL, \yTop+0.45) {Probabilistic Execution};
\node[hdr, text=modelnative!80!black, fill=modelnative!10]
    at (\cxR, \yTop+0.45) {Deterministic Control};

% ============================================================
%  INPUT  (spans both lanes, centred)
% ============================================================
\node[procN, minimum width=13.0cm]
    (req) at (0, 0.75)
    {User request arrives: goal, constraints, environment handle};

% ============================================================
%  LEFT LANE -- Probabilistic Execution  (blue)
% ============================================================
\node[procB] (plan) at (\cxL, 0.0)
    {Agent kernel parses task and loads\\base policies, project memory};

\node[procB, below=4.5mm of plan]
    (perm) {Permission \& capability check\\(deterministic gate)};

\node[procB, below=4.5mm of perm]
    (prefill) {Context compilation \& address resolution:\\
    hot-window loading, summary backfill,\\
    prefix-cache hit, prefill/decode scheduling};

\node[procB, below=4.5mm of prefill]
    (gen) {Model execution \& control loop:\\
    generate action, validate result,\\
    retry or fork if needed};

% ============================================================
%  RIGHT LANE -- Deterministic Control  (orange)
% ============================================================
\node[procO]
    (policy) at (\cxR, 0.0)
    {Load permission config,\\
    tool-capability annotations};

\node[dec, below=6mm of policy,
      draw=classical!50, fill=modelnative!8]
    (tool) {Tool / sub-agent\\needed?};

% Split-colour diagonal lines inside diamond
\draw[modelnative!40, line width=0.4pt] (tool.south west) -- (tool.north east);
\draw[classical!40, line width=0.4pt]   (tool.north west) -- (tool.south east);

\node[procO, below=5mm of tool]
    (syscall) {System call path:\\
    MCP / tool invocation, sandboxed\\
    execution, result return};

\node[procO, below=4.5mm of syscall]
    (commit) {State commit: update short-term memory,\\
    long-term memory, execution log \&\\
    auditable evidence};

% ============================================================
%  CONVERGENCE / OUTPUT  (spans both lanes)
% ============================================================
\node[procN, minimum width=13.0cm]
    (resp) at (0, \yBot)
    {Output response or produce committable artifact};

% ---- Convergence dot ----
\fill[gray!30] (0, \yBot+0.55) circle (3pt);
\node[font=\tiny, text=gray!50, right=3pt] at (0.05, \yBot+0.55) {merge};

% ============================================================
%  ARROWS
% ============================================================
% Input splits into both lanes
\coordinate (split) at (0, 0.15);
\draw[arrG] (req.south) -- (split);
\draw[arrG] (split) -| (plan.north);
\draw[arrG] (split) -| (policy.north);

% Left lane (probabilistic)
\draw[arrB] (plan) -- (perm);
\draw[arrB] (perm) -- (prefill);
\draw[arrB] (prefill) -- (gen);

% Right lane (deterministic)
\draw[arrO] (policy) -- (tool);
\draw[arrO] (tool.south) --
    node[right, font=\scriptsize, text=modelnative!70!black] {Yes}
    (syscall);
\draw[arrO] (syscall) -- (commit);

% Cross-lane: decision "No" feeds context compilation
\draw[arrB] (tool.west) --
    node[above, font=\scriptsize, text=classical!70!black] {No}
    (prefill.east);

% Cross-lane: model generation feeds state commit
\draw[arrO] (gen.east) -- (commit.west);

% Convergence to output
\coordinate (merge) at (0, \yBot+0.9);
\draw[arrG] (gen.south)  |- (merge);
\draw[arrG] (commit.south) |- (merge);
\draw[arrG] (merge) -- (resp.north);

\end{tikzpicture}}
\caption{Request flow through the model-native computing architecture as a dual-plane swim-lane diagram.
Blue nodes belong to the probabilistic execution plane; orange nodes belong to the deterministic control plane.
The split-coloured diamond represents the control-plane decision that gates both paths.
Grey boxes span both planes.}
\label{fig:flow}
\end{figure}

%% ============================================================
